%!
% @file 01_introduction.tex
% @brief Introduction chapter for WordSprint report.

\chapter{Giới thiệu}
\label{chap:introduction}

% -------------------------------------------------
% Tổng quan đề tài
% -------------------------------------------------
\section{Tổng quan đề tài}
\label{sec:overview}
WordSprint là một phần mềm desktop hỗ trợ người học tiếng Anh thông qua các trò chơi đố từ vựng có thời gian giới hạn. Ứng dụng kết hợp các dạng câu đố (xáo trộn chữ, điền từ thiếu, nhận dạng âm thanh) nhằm rèn luyện khả năng nhận diện từ, tốc độ phản ứng và kỹ năng gõ bàn phím. Ý tưởng xuất phát từ nhu cầu tạo môi trường học tập tích cực, vừa thú vị, vừa có tính thử thách, phù hợp cho nhiều trình độ CEFR (A1‑C2) và cho phép người dùng luyện tập theo chủ đề.

% -------------------------------------------------
% Mục tiêu nghiên cứu
% -------------------------------------------------
\section{Mục tiêu nghiên cứu}
\label{sec:objectives}
Mục tiêu chính của dự án WordSprint bao gồm:
\begin{itemize}
    \item \textbf{Cung cấp một công cụ học từ vựng đa dạng} với các loại câu đố được thiết kế để kích thích trí nhớ và phản xạ.
    \item \textbf{Hỗ trợ người dùng ở mọi trình độ} bằng cách cho phép lựa chọn mức độ khó dựa trên tiêu chuẩn CEFR.
    \item \textbf{Xây dựng kiến trúc phần mềm mở rộng được}: giao diện người dùng React + Vite, nền tảng desktop C++/Qt6, cho phép tích hợp thêm nội dung hoặc chức năng mới trong tương lai.
    \item \textbf{Đánh giá hiệu quả học tập} thông qua thống kê thời gian hoàn thành, số câu trả lời đúng và cơ chế lưu trữ hồ sơ người dùng.
\end{itemize}

% -------------------------------------------------
% Phạm vi và đối tượng
% -------------------------------------------------
\section{Phạm vi và đối tượng}
\label{sec:scope}
Hệ thống được thiết kế cho:
\begin{itemize}
    \item \textbf{Người học tiếng Anh} (tự học hoặc trong môi trường lớp học) muốn luyện tập từ vựng một cách nhanh gọn và có tính gamification.
    \item \textbf{Giáo viên} có thể tạo ra các bộ câu hỏi tùy chỉnh và theo dõi tiến độ học của học sinh.
    \item \textbf{Các nhà phát triển} muốn mở rộng hoặc tùy biến giao diện và chức năng nhờ kiến trúc mô-đun.
\end{itemize}

Phạm vi chức năng chính của WordSprint được tóm tắt trong Bảng~\ref{tab:functional_scope}.

\begin{table}[h!]
\centering
\caption{Phạm vi chức năng của hệ thống WordSprint}
\label{tab:functional_scope}
\begin{tabular}{p{0.35\linewidth}p{0.55\linewidth}}
\toprule
Chức năng & Mô tả \\
\midrule
Quản lý trò chơi & Khởi động, tạm dừng, lưu trữ trạng thái và hiển thị điểm số. \\
Câu đố đa dạng & Scramble, Fill‑in‑the‑Blank, Audio với thời gian giới hạn. \\
Hồ sơ người dùng & Đăng ký, chuyển đổi hồ sơ, lưu tiến độ và thành tựu. \\
Bảng xếp hạng & Xếp hạng cá nhân và toàn hệ thống dựa trên điểm số. \\
Cấu hình âm thanh & Điều chỉnh âm lượng âm thanh và nhạc nền. \\
\bottomrule
\end{tabular}
\end{table}

% -------------------------------------------------
% Bố cục báo cáo
% -------------------------------------------------
\section{Bố cục báo cáo}
\label{sec:structure}
Báo cáo này được tổ chức theo các chương sau:
\begin{itemize}
    \item \textbf{Chương 1 -- Giới thiệu} (nội dung hiện tại): cung cấp bối cảnh, mục tiêu, phạm vi và cấu trúc báo cáo.
    \item \textbf{Chương 2 -- Tổng quan tài liệu}: khảo sát các công nghệ liên quan và các nghiên cứu hiện có.
    \item \textbf{Chương 3 -- Kiến trúc hệ thống}: mô tả thiết kế phần mềm, mô‑đun và quy trình giao tiếp.
    \item \textbf{Chương 4 -- Triển khai}: chi tiết quá trình xây dựng, biên dịch và tích hợp giao diện.
    \item \textbf{Chương 5 -- Kết luận và hướng phát triển}: tổng kết kết quả, đánh giá và đề xuất cải tiến trong tương lai.
\end{itemize}

\noindent Nội dung các chương được liên kết chặt chẽ qua các nhãn (\ref{chap:introduction}, \ref{sec:overview}, \ref{sec:objectives}, \ref{sec:scope}, \ref{sec:structure}) nhằm hỗ trợ người đọc theo dõi một cách hệ thống.
